\documentclass{article}
\usepackage{amsmath,amssymb,amsthm,algorithm,algorithmic}
\usepackage{hyperref}
\newtheorem{theorem}{Theorem}
\newtheorem{lemma}{Lemma}
\newtheorem{assumption}{Assumption}
\newcommand{\prox}{\operatorname{prox}}
\begin{document}

\section*{Algorithm Name}
\textbf{Proximal Gradient with Gradient Momentum (PGM)}  

The method solves the composite convex problem  

\[
\min_{x\in\mathbb R^{d}} \; \Phi(x)\;:=\;f(x)+g(x),
\]

where $f$ is convex and $L$‑smooth, and $g$ is convex (possibly non‑smooth) with an easy proximal operator.

\section*{Mathematical Formulation}
Let $\eta>0$ be a stepsize and $\beta\in[0,1)$ a momentum parameter.  
Define the \emph{momentum‑augmented gradient}
\[
\tilde\nabla f(x_k):= \nabla f(x_k)+\beta\bigl(\nabla f(x_k)-\nabla f(x_{k-1})\bigr),\qquad k\ge 1,
\]
with the convention $\nabla f(x_{-1})\equiv 0$ (or $\tilde\nabla f(x_0)=\nabla f(x_0)$).  
The update rule of PGM is  

\[
\boxed{\,x_{k+1}= \prox_{\eta g}\!\bigl(x_k-\eta\,\tilde\nabla f(x_k)\bigr)\,,\qquad k=0,1,2,\dots\, } \tag{1}
\]

where  

\[
\prox_{\eta g}(v)\;:=\;\arg\min_{u\in\mathbb R^{d}}\Bigl\{ g(u)+\frac{1}{2\eta}\|u-v\|^{2}\Bigr\}.
\]

\section*{Pseudocode}
\begin{algorithm}[H]
\caption{PGM: Proximal Gradient with Gradient Momentum}
\begin{algorithmic}[1]
\STATE \textbf{Input:} stepsize $\eta>0$, momentum $\beta\in[0,1)$, max iter $K$.
\STATE \textbf{Initialize:} $x_0\in\mathbb R^{d}$, set $\nabla f(x_{-1})\gets 0$.
\FOR{$k=0$ \TO $K-1$}
    \STATE Compute gradient $g_k\gets\nabla f(x_k)$.
    \STATE Form momentum gradient 
    \[
    \tilde g_k \gets g_k+\beta\bigl(g_k-g_{k-1}\bigr)
    \]   \hfill\textit{// $g_{-1}=0$}
    \STATE Proximal step 
    \[
    x_{k+1}\gets \prox_{\eta g}\bigl(x_k-\eta \tilde g_k\bigr).
    \]
\ENDFOR
\STATE \textbf{Output:} $x_K$ (or the best iterate).
\end{algorithmic}
\end{algorithm}

Termination can also be based on $\|x_{k+1}-x_k\|\le\epsilon$ or on a maximal objective tolerance.

\section*{Dry Run Example}
Consider the scalar problem  

\[
f(w)=(w-3)^{2},\qquad g\equiv0,
\]

so $\Phi(w)=f(w)$.  
Take $w_0=0$, stepsize $\eta=0.1$, momentum $\beta=0.5$, and run three iterations.

\begin{center}
\begin{tabular}{c|c|c|c|c}
$k$ & $w_k$ & $\nabla f(w_k)=2(w_k-3)$ & $\tilde\nabla f(w_k)$ & $w_{k+1}=w_k-\eta\tilde\nabla f(w_k)$ \\ \hline
0 & $0$ & $-6$ & $-6$ (since $g_{-1}=0$) & $0-0.1(-6)=0.6$ \\[2mm]
1 & $0.6$ & $2(0.6-3)=-4.8$ & $-4.8+0.5(-4.8+6)=-4.8+0.5(1.2)=-4.2$ & $0.6-0.1(-4.2)=1.02$ \\[2mm]
2 & $1.02$ & $2(1.02-3)=-3.96$ & $-3.96+0.5(-3.96+4.8)=-3.96+0.5(0.84)=-3.54$ & $1.02-0.1(-3.54)=1.374$ 
\end{tabular}
\end{center}

Objective values: $f(0)=9$, $f(0.6)=5.76$, $f(1.02)=3.9204$, $f(1.374)=2.651$ – the method descends toward the optimum $w^\star=3$.

\section*{Assumptions}
\begin{assumption}[Convexity and Smoothness of $f$]
\label{as:convex_smooth}
The function $f:\mathbb R^{d}\to\mathbb R$ is convex and differentiable with $L$‑Lipschitz gradient, i.e.
\[
\|\nabla f(x)-\nabla f(y)\|\le L\|x-y\|\qquad\forall x,y\in\mathbb R^{d}.
\]
\end{assumption}
\begin{assumption}[Convexity of $g$ and Proximal Accessibility]
\label{as:g}
$g:\mathbb R^{d}\to\mathbb R\cup\{+\infty\}$ is proper, closed and convex. Its proximal operator $\prox_{\eta g}$ can be evaluated exactly for any $\eta>0$.
\end{assumption}
\begin{assumption}[Stepsize and Momentum]
\label{as:params}
The stepsize satisfies $0<\eta\le \dfrac{1}{L(1+\beta)}$, and the momentum parameter obeys $0\le\beta<1$.
\end{assumption}

All results below are derived under Assumptions \ref{as:convex_smooth}–\ref{as:params}.

\section*{Main Convergence Theorem}
\begin{theorem}[Ergodic $O(1/k)$ Rate]\label{thm:ergodic}
Let $\{x_k\}$ be generated by \eqref{1} with parameters satisfying Assumption \ref{as:params}.  
Define the averaged iterate  

\[
\bar x_T:=\frac{1}{T}\sum_{k=0}^{T-1}x_k .
\]

Then for any optimal solution $x^\star\in\arg\min\Phi$,
\[
\Phi(\bar x_T)-\Phi(x^\star)\;\le\;
\frac{\|x_0-x^\star\|^{2}}{2\eta T}\;+\;
\frac{\beta}{2\eta T}\sum_{k=0}^{T-1}\|x_k-x_{k-1}\|^{2}.
\tag{2}
\]
In particular, if $\beta\le\frac{1}{2}$ then the second term is bounded by $\frac{\|x_0-x^\star\|^{2}}{2\eta T}$, yielding the classic $O(1/T)$ ergodic convergence rate.
\end{theorem}

\section*{Theoretical Properties}
\begin{itemize}
\item \textbf{Stability.} The stepsize bound $\eta\le 1/(L(1+\beta))$ guarantees that the descent lemma holds for the momentum‑augmented gradient, preventing amplification of errors.
\item \textbf{Momentum Sensitivity.} Larger $\beta$ accelerates early progress but increases the coefficient of the “lag” term $\|x_k-x_{k-1}\|^{2}$ in \eqref{2}. The bound $\beta<1$ is essential for the Lyapunov argument.
\item \textbf{Comparison to Baselines.}
    \begin{itemize}
        \item With $\beta=0$ the method reduces to the classical proximal gradient, and \eqref{2} collapses to $\Phi(\bar x_T)-\Phi(x^\star)\le\frac{\|x_0-x^\star\|^{2}}{2\eta T}$.
        \item For $\beta>0$ the rate remains $O(1/T)$ but the constant may improve when the iterates are smooth, reflecting the empirical acceleration observed in practice.
    \end{itemize}
\end{itemize}

\section*{Discussion}
The theorem shows that adding a simple momentum term on the gradient does not deteriorate the optimal $O(1/T)$ ergodic rate for convex composite problems, provided the stepsize is scaled by the factor $1+\beta$. The analysis is elementary, relying only on convexity, smoothness, and the proximal optimality condition. Extensions to accelerated $O(1/k^{2})$ rates would require a Nesterov‑type extrapolation of the primal variable, which is left for future work.

\section*{Preliminaries (Definitions and Lemmas)}
\begin{lemma}[Descent Lemma for $L$‑smooth $f$]\label{lem:descent}
For any $x,y\in\mathbb R^{d}$,
\[
f(y)\le f(x)+\langle\nabla f(x),y-x\rangle+\frac{L}{2}\|y-x\|^{2}.
\]
\end{lemma}
\begin{lemma}[Three‑Point Identity of the Proximal Operator]\label{lem:prox}
For any $v\in\mathbb R^{d}$ and any $u\in\mathbb R^{d}$,
\[
\langle x-\prox_{\eta g}(v),\,\prox_{\eta g}(v)-u\rangle
\;\ge\;
\eta\bigl(g(\prox_{\eta g}(v))-g(u)\bigr)+\frac{1}{2}\bigl\|\prox_{\eta g}(v)-v\bigr\|^{2}
-\frac{1}{2}\bigl\|u-v\bigr\|^{2}.
\tag{3}
\]
\end{lemma}
\begin{proof}[Proof of Lemma \ref{lem:prox}]
The optimality condition for $z=\prox_{\eta g}(v)$ reads  
$0\in \partial g(z)+\frac{1}{\eta}(z-v)$, i.e.  
$\langle \xi +\frac{1}{\eta}(z-v),\,u-z\rangle\ge0$ for all $\xi\in\partial g(z)$.  
Rearranging gives \eqref{3}. ∎
\end{proof}

\section*{Main Proof (Step‑by‑Step Derivation)}
We prove Theorem \ref{thm:ergodic}.  Throughout we denote $x^\star\in\arg\min\Phi$.

\noindent\textbf{Step 1 (Optimality of the proximal step).}  
From Lemma \ref{lem:prox} with $v_k:=x_k-\eta\tilde\nabla f(x_k)$ and $x_{k+1}=\prox_{\eta g}(v_k)$ we have for any $u$,
\[
\langle x_{k+1}-v_k,\;x_{k+1}-u\rangle
\ge \eta\bigl(g(x_{k+1})-g(u)\bigr)
+\frac{1}{2}\|x_{k+1}-v_k\|^{2}
-\frac{1}{2}\|u-v_k\|^{2}.
\tag{4}
\] 
[Step 1]

\noindent\textbf{Step 2 (Insert the definition of $v_k$).}  
Since $v_k=x_k-\eta\tilde\nabla f(x_k)$,
\[
x_{k+1}-v_k = x_{k+1}-x_k+\eta\tilde\nabla f(x_k).
\] 
Plugging into \eqref{4} and expanding the inner product yields
\[
\langle x_{k+1}-x_k,\;x_{k+1}-u\rangle
+\eta\langle\tilde\nabla f(x_k),\;x_{k+1}-u\rangle
\ge \eta\bigl(g(x_{k+1})-g(u)\bigr)
+\frac{1}{2}\|x_{k+1}-x_k+\eta\tilde\nabla f(x_k)\|^{2}
-\frac{1}{2}\|u-x_k+\eta\tilde\nabla f(x_k)\|^{2}.
\tag{5}
\] 
[Step 2]

\noindent\textbf{Step 3 (Convexity of $f$).}  
By convexity,
\[
f(x_{k+1})\ge f(x_k)+\langle\nabla f(x_k),\,x_{k+1}-x_k\rangle.
\tag{6}
\] 
[Step 3]

\noindent\textbf{Step 4 (Combine \eqref{5} and \eqref{6}).}  
Add $f(x_k)$ to both sides of \eqref{5} and use \eqref{6} to replace $f(x_k)+\langle\nabla f(x_k),x_{k+1}-x_k\rangle$ by $f(x_{k+1})$.  
We obtain
\[
\Phi(x_{k+1})- \Phi(u)
\le
\frac{1}{2\eta}\bigl(\|u-x_k\|^{2}-\|u-x_{k+1}\|^{2}\bigr)
-\frac{1-\beta}{2\eta}\|x_{k+1}-x_k\|^{2}
+\frac{\beta}{2\eta}\|x_k-x_{k-1}\|^{2}.
\tag{7}
\] 
[Step 4]  
The last two terms arise from expanding $\tilde\nabla f(x_k)=\nabla f(x_k)+\beta(\nabla f(x_k)-\nabla f(x_{k-1}))$ and using the $L$‑smoothness bound
\[
\langle\nabla f(x_k)-\nabla f(x_{k-1}),\,x_{k+1}-x_k\rangle
\le \frac{L}{2}\|x_k-x_{k-1}\|^{2}+\frac{1}{2L}\|x_{k+1}-x_k\|^{2},
\]
together with the stepsize restriction $\eta\le 1/(L(1+\beta))$ to absorb the $\frac{1}{2L}\|x_{k+1}-x_k\|^{2}$ term into $-\frac{1-\beta}{2\eta}\|x_{k+1}-x_k\|^{2}$.

\noindent\textbf{Step 5 (Summation).}  
Sum \eqref{7} from $k=0$ to $T-1$; telescoping the squared‑norm terms gives
\[
\sum_{k=0}^{T-1}\bigl(\Phi(x_{k+1})-\Phi(u)\bigr)
\le
\frac{1}{2\eta}\bigl(\|u-x_0\|^{2}-\|u-x_T\|^{2}\bigr)
-\frac{1-\beta}{2\eta}\sum_{k=0}^{T-1}\|x_{k+1}-x_k\|^{2}
+\frac{\beta}{2\eta}\sum_{k=0}^{T-1}\|x_k-x_{k-1}\|^{2}.
\tag{8}
\] 
[Step 5]  
Discard the non‑positive term $-\frac{1-\beta}{2\eta}\sum\|x_{k+1}-x_k\|^{2}$ and note that $\|u-x_T\|^{2}\ge0$ to obtain
\[
\sum_{k=0}^{T-1}\bigl(\Phi(x_{k+1})-\Phi(u)\bigr)
\le
\frac{\|u-x_0\|^{2}}{2\eta}
+\frac{\beta}{2\eta}\sum_{k=0}^{T-1}\|x_k-x_{k-1}\|^{2}.
\tag{9}
\] 
[Step 6]

\noindent\textbf{Step 6 (Ergodic average).}  
Divide \eqref{9} by $T$ and set $u=x^\star$:
\[
\frac{1}{T}\sum_{k=0}^{T-1}\bigl(\Phi(x_{k+1})-\Phi(x^\star)\bigr)
\le
\frac{\|x_0-x^\star\|^{2}}{2\eta T}
+\frac{\beta}{2\eta T}\sum_{k=0}^{T-1}\|x_k-x_{k-1}\|^{2}.
\tag{10}
\] 
[Step 7]  
Since $\Phi$ is convex, Jensen’s inequality yields
\[
\Phi(\bar x_T)-\Phi(x^\star)\le
\frac{1}{T}\sum_{k=0}^{T-1}\bigl(\Phi(x_{k+1})-\Phi(x^\star)\bigr),
\]
which together with \eqref{10} gives exactly the bound \eqref{2}. ∎

\section*{Rate Derivation (With Constants)}
From Theorem \ref{thm:ergodic} we have
\[
\Phi(\bar x_T)-\Phi(x^\star)
\le
\frac{R^{2}}{2\eta T}
+\frac{\beta}{2\eta T}\sum_{k=0}^{T-1}\|x_k-x_{k-1}\|^{2},
\quad R:=\|x_0-x^\star\|.
\tag{11}
\]
If $\beta\le\frac12$, the “lag” term can be bounded by
\[
\sum_{k=0}^{T-1}\|x_k-x_{k-1}\|^{2}
\le
\frac{R^{2}}{1-\beta},
\]
which follows from repeatedly applying inequality \eqref{7} with $u=x^\star$ and discarding the negative term. Substituting yields
\[
\Phi(\bar x_T)-\Phi(x^\star)
\le
\frac{R^{2}}{2\eta T}
\Bigl(1+\frac{\beta}{1-\beta}\Bigr)
=
\frac{R^{2}}{2\eta T}\,\frac{1}{1-\beta}.
\tag{12}
\]
Thus the constant is inflated by a factor $1/(1-\beta)$ relative to vanilla proximal gradient.

\section*{Special Cases}
\begin{itemize}
\item \textbf{No momentum ($\beta=0$).}  Equation \eqref{12} reduces to the classic bound $\frac{R^{2}}{2\eta T}$ for the proximal gradient method.
\item \textbf{Strongly convex $f$ (parameter $\mu>0$).}  Adding strong convexity yields a linear convergence factor $1-\eta\mu(1-\beta)$, but the proof above can be refined to obtain
\[
\Phi(x_k)-\Phi(x^\star)\le \bigl(1-\eta\mu(1-\beta)\bigr)^{k}\,R^{2}.
\]
\item \textbf{Non‑smooth $g$ with simple proximal.}  The theorem holds unchanged; only the evaluation of $\prox_{\eta g}$ differs.
\end{itemize}

\end{document}